\documentclass[10pt]{article}
\usepackage[margin=1in]{geometry}
\usepackage[T1]{fontenc}
\usepackage{lmodern}
\usepackage{amsmath,amssymb}
\usepackage{booktabs}
\usepackage{graphicx}
\usepackage{xcolor}
\usepackage[numbers,sort&compress]{natbib}
\usepackage[hidelinks]{hyperref}

% Every number in this paper comes from paper/generated/*.tex or results/figures/,
% both regenerated from the experiment DB by ./reproduce.sh. Missing results render
% as a visible PENDING box; nothing is typed by hand.
\newcommand{\pending}[1]{\fbox{\textcolor{red}{\textsc{pending}}: \small #1}}
\newcommand{\fig}[3]{%
  \IfFileExists{#1}{\includegraphics[width=#2\linewidth]{#1}}{\pending{#3}}}

\title{Is Compressibility Predictable Early?\\
\large Compression-aware training, a measured lifecycle, and training-time signals}
\author{Eshaan Kapooswalla}
\date{}

\begin{document}
\maketitle

\begin{abstract}
We ask whether compression-awareness during training yields models that quantize and prune
better after training and run faster at inference, and whether a model's final
compressibility can be predicted early in training from cheap signals. We build a
measured lifecycle: compression-aware training variants (Quant-Noise, kurtosis
regularization, RigL), a post-training ladder (RTN with calibrated activations, AdaRound,
BRECQ, LSQ fine-tuning, Hessian-aware mixed precision, SmoothQuant, N:M and channel
pruning with masked LoRA recovery), single-threaded CPU kernels (fp32, exact int8, int4
weight-only, 2:4 and CSR sparse), and a benchmarking protocol that treats fresh processes as
the unit of replication. Signals logged during training (weight kurtosis, weight and
activation outlier ratios, Hessian trace, sharpness) are features for predicting the final
model's accuracy drop under a fixed compression panel, evaluated on held-out training
configurations. \textbf{Headline finding:} \pending{written after the Phase 5/6 runs; see
Section~\ref{sec:results}}.
\end{abstract}

\section{Introduction}
Model compression is usually studied one stage at a time: a quantization method is evaluated
on a fixed pre-trained checkpoint, and a pruning method on another. Yet the checkpoint's
training history plausibly determines how much compression it tolerates. Weight
distributions with heavy tails waste quantization levels~\citep{shkolnik2020robust}, and
flat minima tolerate weight perturbations~\citep{keskar2017large,foret2021sam}. If
compressibility is a property of training, then (i) training can be made compression-aware,
and (ii) compressibility may be \emph{predictable} before training ends, which would let
practitioners discard poor candidates early.

We study both questions on image classification at CIFAR scale, with every measurement
produced by one reproducible pipeline:
\begin{itemize}
  \item \textbf{A measured lifecycle.} Training variants, a post-training ladder, lowering to
  our own CPU kernels, and latency/memory benchmarking. All results live in one experiment
  database, and every reported number traces to a run identifier.
  \item \textbf{A benchmarking protocol} that uses fresh processes as the unit of
  replication~\citep{kalibera2013rigorous} and records a hardware/software fingerprint for
  every run.
  \item \textbf{An early-predictability study} with leave-one-configuration-out evaluation,
  baselines that control for model size, and signal ablations.
  \item \textbf{A latency-proxy study} testing whether FLOPs, parameter count or sparsity
  predict measured latency across kernel families, compared with a learned per-layer proxy.
\end{itemize}

\section{Related work}
\paragraph{Post-training quantization.} Rounding to nearest with calibrated activation
ranges is the standard baseline~\citep{krishnamoorthi2018,jacob2018quantization,nagel2021white}.
AdaRound learns per-weight rounding by layer-wise output reconstruction~\citep{nagel2020adaround};
BRECQ reconstructs whole blocks~\citep{li2021brecq}; LSQ learns step sizes during
fine-tuning~\citep{esser2020lsq}; HAWQ allocates bits by Hessian
sensitivity~\citep{dong2020hawqv2}; SmoothQuant migrates activation outliers into
weights~\citep{xiao2023smoothquant}, motivated by outliers that emerge at large
scale~\citep{dettmers2022llmint8}.
\paragraph{Compression-aware training.} Quant-Noise quantizes random weight subsets during
training~\citep{fan2021quantnoise}; kurtosis regularization makes weights uniform-like and
robust to quantization~\citep{shkolnik2020robust}; RigL trains sparse networks by dynamic
drop-and-grow~\citep{evci2020rigl}.
\paragraph{Pruning.} Structured N:M sparsity targets hardware
support~\citep{mishra2021accelerating,zhou2021nm}; channel pruning ranks filters by norm or
BN scale~\citep{li2017pruning,liu2017slimming}; low-rank adapters~\citep{hu2022lora} are a
parameter-efficient recovery option.
\paragraph{Early signals.} Early-Bird tickets predict pruning masks early in
training~\citep{you2020earlybird}; curvature and sharpness measures relate to
generalization~\citep{yao2020pyhessian,jiang2020fantastic}, with known caveats about
reparametrization~\citep{dinh2017sharp}.
\paragraph{Latency.} FLOPs is an indirect proxy for speed~\citep{ma2018shufflenetv2}; learned
latency predictors are an alternative~\citep{zhang2021nnmeter}. We use the roofline
model~\citep{williams2009roofline} to explain kernel behaviour.

\section{System}
\paragraph{Benchmarking.} Each measurement runs in $R$ fresh Python processes. Each process
records its cold-start stages, warms up, times at least $N$ calls with the garbage collector
disabled, and reports its peak resident memory (\texttt{/proc} VmHWM on Linux; see
Appendix~\ref{app:pitfalls}). We report the median of per-process medians with a
percentile-bootstrap CI over processes, and pooled percentiles only descriptively, because
iterations within a process are autocorrelated~\citep{kalibera2013rigorous,mytkowicz2009wrong}.
\paragraph{Kernels and engine.} Single-threaded C++ GEMMs with runtime ISA dispatch (NEON with
dot-product instructions; AVX2/AVX-512): fp32 with packed weight panels, exact
int8$\times$int8$\rightarrow$int32, int4 weight-only with per-group scales, 2:4 structured
and CSR sparse. A torch.fx-based engine folds BatchNorm, fuses ReLU and lowers convolutions to
im2col+GEMM. For quantized models it uses exactly the integer weights and scales of the
simulated model. Every lowered layer matches its simulation to $10^{-5}$ relative given
identical inputs (tested). End to end the two agree up to rounding-boundary flips.
\paragraph{Experiment database.} One row per run with configuration tags, accuracy, latency
statistics, memory, measured weight sparsity, fingerprint hash and git commit, plus per-run
artifacts. Figures and every table in this paper are regenerated from the database by
\texttt{reproduce.sh}.

\section{Method}
\paragraph{Training variants.} Standard SGD; Quant-Noise (4-bit, $p{=}0.5$) and its $p{=}1$
limit, i.e.\ quantization-aware training from scratch; kurtosis regularization toward 1.8; RigL at
90\% sparsity with ERK allocation.
\paragraph{Signals.} At initialization, at log-spaced checkpoints and at a fixed interval, on
a fixed probe batch of \emph{training} data: per-layer weight kurtosis, max$|w|$/std, weight
norms, input-activation outlier ratios, the Hutchinson trace of the loss
Hessian~\citep{hutchinson1989}, and first-order SAM sharpness~\citep{foret2021sam}.
\paragraph{Compressibility target.} The accuracy drop of the final model under a fixed,
training-free panel: RTN W8A8, RTN W4A8, W4 weight-only (group 32), 2:4, and 50\%/80\%
unstructured.
\paragraph{Prediction protocol.} Features are the signals at training fraction $f$, plus
$\log_{10}$ parameter count and a variant one-hot. Predictors (random forest, gradient
boosting, MLP, Gaussian process) are compared with a mean baseline and a parameters-only
baseline. Evaluation holds out every seed of one configuration at a time. CIs come from a
cluster bootstrap over configurations.

\section{Experimental setup}
CIFAR-10~\citep{krizhevsky2009cifar}. ResNets~\citep{he2016resnet} of depth 10/18/34 and width
0.25/0.5/1.0, and three ViTs~\citep{dosovitskiy2021vit} (patch 4). Three training seeds per
configuration. Latency is measured on an Apple M5 Pro (primary platform, OS-scheduled, all
backends at one thread) and on x86 cloud VMs (secondary, noisy). The CIs account for the
noise, and the platform caveats are listed in Section~\ref{sec:limits}.

\section{Results}\label{sec:results}
\subsection{The post-training ladder}
\begin{center}\small\pending{run configs/sweeps/ptq\_ladder\_resnet18\_cifar10.yml}
\end{center}
\subsection{Pruning and recovery}
\begin{center}\small\pending{run configs/sweeps/prune\_ladder\_resnet18\_cifar10.yml}
\end{center}
\subsection{Early predictability of compressibility}
\begin{center}
\fig{../results/figures/early_prediction_w4a8.png}{0.7}{run the Phase 5 tracks + compress\_runs, then reproduce.sh}\\[4pt]
\small\pending{run the Phase 5 tracks + compress\_runs, then reproduce.sh (w4a8)}

\end{center}
\subsection{Scaling}
\begin{center}\small\pending{run pretrain\_scaling + compress\_runs, then reproduce.sh}
\end{center}
\subsection{Do FLOPs predict latency?}
\begin{center}\small\pending{run the kernel study and backend\_latency\_m5 on the M5 Pro, then reproduce.sh}
\end{center}
\subsection{When does sparsity pay on a CPU?}
\begin{center}
\fig{../results/figures/kernel_sparsity_m5_sparsity.png}{0.6}{run configs/studies/kernel\_sparsity\_m5.yml}\\[4pt]
\small\pending{run configs/studies/kernel\_sparsity\_m5.yml, then reproduce.sh}

\end{center}

\subsection{Stretch: small language models}
GPTs of roughly 10M/25M/50M parameters (byte-level tokens) are pre-trained on
TinyStories~\citep{eldan2023tinystories},
instruction-tuned on TinyStories-Instruct, and aligned with DPO~\citep{rafailov2023dpo} toward
satisfying the instruction's required-words constraint. The same quantization ladder is then
applied to each stage. \emph{Capability} is held-out perplexity; \emph{behavior} is the fraction
of greedy completions containing every required word. The question is whether behavior
degrades before capability. Speculative decoding~\citep{leviathan2023speculative} with a
distilled draft is lossless for greedy decoding; we report its acceptance rate.
\begin{center}\small\pending{run configs/lm/tinystories\_gpt\_*.yml (3 seeds), then reproduce.sh}
\end{center}

\section{Limitations}\label{sec:limits}
\begin{itemize}
  \item CIFAR-scale models and data. ImageNet validation against the published PTQ numbers
  needs license-gated data and is \pending{run configs/sweeps/ptq\_validation\_imagenet.yml}.
  \item Kernels are single-threaded, so absolute latencies understate what a multi-threaded
  runtime achieves.
  \item Latency is measured on one primary machine, which cannot pin threads to cores. Energy
  is not measured.
  \item Three seeds per configuration. CIs are percentile bootstraps and under-cover at small $n$.
  \item Activations are symmetric int8, including after ReLU (one bit unused), to match the
  kernels.
  \item The prediction targets use a training-free panel. Methods with recovery or
  reconstruction may change the ranking.
\end{itemize}

\section*{Reproducibility}
Every experiment is one configuration file and one command. \texttt{./reproduce.sh}
regenerates all figures, and all tables in this paper, from the experiment database. Design
decisions, including the pitfalls below, are logged in \texttt{docs/DECISIONS.md}.

\appendix
\section{Measurement pitfalls found while building the pipeline}\label{app:pitfalls}
Each of these would have silently corrupted results:
\begin{itemize}
  \item On Linux, \texttt{ru\_maxrss} in a spawned worker reports the \emph{parent's} peak
  memory: the kernel folds the pre-exec address space's high-water mark in at
  \texttt{execve}. Use \texttt{/proc} VmHWM.
  \item Multiplying by $1/s$ instead of dividing by $s$ when quantizing flips rounding ties,
  so the engine and the simulation diverge.
  \item FLOP counters that match only float layer types report zero FLOPs for quantized models.
  \item A missing checkpoint path silently evaluated randomly initialized weights.
\end{itemize}

\bibliographystyle{plainnat}
\bibliography{refs}
\end{document}
